% !TeX program = xelatex
% Evidence cut: the supplied seed-42 core group, plus the completed 8-crop diagnostic.
% Ours = search-context C architecture + one LR decay after update 936.
\documentclass[journal]{IEEEtran}
\usepackage{fontspec}
\IfFontExistsTF{TeX Gyre Termes}{\setmainfont{TeX Gyre Termes}}{\setmainfont{Latin Modern Roman}}
\usepackage{xeCJK}
\IfFontExistsTF{SimSun}{%
  \IfFontExistsTF{SimHei}{\setCJKmainfont[BoldFont=SimHei]{SimSun}}{\setCJKmainfont[AutoFakeBold=true]{SimSun}}%
}{\setCJKmainfont[AutoFakeBold=true]{FandolSong-Regular}}
\IfFontExistsTF{Microsoft YaHei}{\setCJKsansfont{Microsoft YaHei}}{\setCJKsansfont{FandolHei-Regular}}
\IfFontExistsTF{SimSun}{\setCJKmonofont{SimSun}}{\setCJKmonofont{FandolFang-Regular}}
\usepackage{cite}
\usepackage{amsmath,amssymb,amsfonts}
\usepackage{graphicx,booktabs,array}
\usepackage{balance}
\usepackage[hidelinks]{hyperref}
\renewcommand{\abstractname}{摘\quad 要}
\renewcommand{\figurename}{图}
\renewcommand{\tablename}{表}
\renewcommand{\IEEEkeywordsname}{关键词}
\newcommand{\R}{\mathrm{R}}
\newcommand{\T}{\mathrm{T}}
\newcommand{\softmax}{\operatorname{softmax}}
\newcommand{\LN}{\operatorname{LN}}
\newcommand{\clip}{\operatorname{clip}}
\newcommand{\NA}{\mathrm{NA}}
\begin{document}

\title{搜索上下文条件化模板 Value 交互的 RGB-T 目标跟踪}
\author{\textit{作者与单位待补充}}
\maketitle

\begin{abstract}
RGB-T 单目标跟踪利用可见光与热红外信息持续定位目标。
模板跨模态注意力能够建立内容匹配，但匹配相关性不等于来源消息的定位效用。
已有 Early gate 将模板与搜索区域的关系压缩为少量标量，再调节跨模态 Value 幅度。
本文进一步研究：显式检索当前搜索内容，能否为模板来源的消息调制提供更充分的条件。
为此，本文构建搜索上下文条件化模板 Value 交互。
来源模板 token 检索对侧当前搜索 token，并结合原有 early evidence 生成残差增益。
该增益仅作用于跨模态来源 Value，保留同模态消息及原有定位路径。
在 LasHeR 全部 244 条测试序列、219,292 帧上，当前候选取得 AUC 56.8305、PR 70.4502、NPR 66.8347、OP50 68.8178、OP75 49.0406。
相对基础系统 B0，五项历史点估计分别提高 2.01、1.92、2.24、2.30、3.24 个百分点。
相对模板 MHSA 参照 M1，AUC 和 OP75 分别提高 0.62 和 1.91 个百分点。
这些历史比较存在训练与来源差异；同结构调度对照及固定裁剪干预另行报告。
结果提供了搜索条件化模板交互的经验性能证据，尚不足以确认退化识别或错误传播纠正机制。
\end{abstract}

\begin{IEEEkeywords}
RGB-T 单目标跟踪，模板交互，搜索上下文，Value 调制，逐序列分析。
\end{IEEEkeywords}

\section{引言}
单目标跟踪根据首帧目标框，在后续视频中持续估计目标位置。
可见光提供纹理和颜色，热红外提供不同成像条件下的目标线索。
照明变化、热交叉、遮挡及背景干扰会改变两种模态的有效信息。
LasHeR 等基准为多场景 RGB-T 跟踪提供了系统评测条件\cite{li2022lasher}。
如何根据当前搜索区域利用两种模态，仍是双模态交互的重要问题。

模板桥接方法利用目标模板组织双模态搜索特征的交互\cite{hui2023tbsi,li2025tbsi}。
在本文研究的模板多头自注意力（MHSA）中，两模态模板共同参与 QK 匹配和 Value 聚合。
QK 分数描述内容相关性，却不直接衡量来源 Value 对最终定位的边际贡献。
模板 token 也可能包含目标边界和背景信息。
因此，高匹配响应不能单独作为有益跨模态消息的判据。

已有 Early gate 引入模板到同模态搜索的峰值相似度，以及对应模板 token 的跨模态一致性。
它利用三个标量，为每个来源 token 生成两个方向的 Value 增益。
这一机制已经依赖当前搜索证据，不能视为完全无条件融合。
但是，取最大值会压缩匹配分布，也不保留被检索搜索位置的内容表示。
对侧搜索区域中与来源模板相关的上下文，因而没有通过显式内容检索进入增益预测器。

本文据此提出一个范围明确的问题：
\emph{在保留模板注意力及既有 Early gate 的基础上，显式检索对侧当前搜索内容，能否改善跨模态来源 Value 的调制及跟踪表现？}
我们以来源模板 token 为查询，以对侧搜索 token 和原有证据摘要为记忆，预测残差增益。
增益沿来源 token 轴作用于跨模态 Value，同模态消息保持原样。
该结构称为搜索上下文条件化 Value 交互，实验中记作 C。
最终报告的 Ours 采用 C 结构，并单独标明提前降低学习率的训练设置。

问题动机还来自跟踪中的系统级失败链：crop 内定位失配可能经预测框反馈放大，随后目标出界且难以恢复。
这一现象没有确定模板 MHSA 为首因，也不能由最终 AUC 反推出某条消息有害。
因此，本文将整体性能、调度贡献和学习后动作效用分别检验。
固定状态下的消息干预与自由跟踪结果承担不同的验证任务。

本文的工作包括：
\begin{itemize}
  \item 将模板交互问题收窄为当前搜索内容对来源 Value 调制的条件作用，区分匹配相关性与定位效用。
  \item 在 Early gate 上引入双向搜索内容检索，以残差 logit 调节跨模态来源 Value，不增加前景输入或额外定位分支。
  \item 在 LasHeR full244 上报告基础系统、模板 MHSA 参照及当前候选，并通过调度配对、短预算对照、尾部分析与固定 crop 干预界定结论。
\end{itemize}

\section{相关工作与问题定位}
\subsection{RGB-T 跟踪与模板交互}
LasHeR 扩展了 RGB-T 跟踪的序列规模与场景覆盖\cite{li2022lasher}。
TBSI 通过模板桥接组织双模态搜索区域交互\cite{hui2023tbsi}。
后续工作进一步研究模板桥接及目标保持的模板更新\cite{li2025tbsi}。
本文以这一研究背景为起点，将修改位置限定在模板跨模态 Value 支路。
本文使用的 B0 是本研究记录中的基础系统参照，不将其分数表述为原论文的统一官方结果。
M1 则表示本研究已有的模板联合 MHSA 实现。

\subsection{内容匹配与消息调制}
注意力权重决定接收 token 如何聚合来源信息，来源增益则决定某个 Value 以何种幅度进入跨模态消息。
两者作用于不同环节。
Early gate 已用搜索相关摘要调制来源幅度；C 进一步保留通过检索获得的搜索内容。
其区别在于增益预测器的条件表示，而非重新发明 cross-attention 或门控。
本文不将来源增益解释为经过校准的模态可靠性概率，也不据算子名称主张广泛新颖性。

\subsection{研究范围}
本文的 C 不输入显式前景 mask，不引入双向预算守恒，也不修正单个 Value 向量的方向。
两个来源方向的增益分别预测，无和为常数的约束。
修改保留当前层原有 QK 计算，但前层 Value 更新可以影响后续特征及注意力。
同一来源增益由所有接收模板 token 和注意力头共享，尚未形成接收端逐 token 的效用控制。
本文同时保留原有定位头、裁剪规则及状态更新方式。

\section{方法}
\subsection{整体计算流程}
设模态 $m\in\{\R,\T\}$，$\bar m$ 为另一模态。
模板和搜索 token 分别记为 $Z_m\in\mathbb{R}^{M\times d}$、$X_m\in\mathbb{R}^{N\times d}$。
当前配置中 $M=64$、$N=256$、$d=768$。
骨干第 3 个 Transformer block 输出的四路特征记为 $Z_m^e,X_m^e$。
它们被保留为本帧 early context，供后续三个融合位置使用。
融合位于零起始 block 索引 3、6、9，即第 4、7、10 个 block 之后。

\begin{figure*}[t]
\centering
\fbox{\begin{minipage}{0.94\textwidth}
\centering
双模态模板与当前搜索图像 $\longrightarrow$ 共享骨干早期特征 $\{Z_\R^e,Z_\T^e,X_\R^e,X_\T^e\}$
\par\medskip
$\underbrace{\text{峰值匹配与模板一致性}}_{e_j}$
$\longrightarrow$ 原有预测器 $f_\ell(e_j)$
\qquad
$\underbrace{\text{来源模板查询对侧搜索及证据摘要}}_{q_j,c_j}$
$\longrightarrow$ 残差预测器 $h_\ell(q_j,c_j,e_j)$
\par\medskip
两路 logit 相加 $\longrightarrow$ 有界来源增益 $g_{\ell,j}^{m\to\bar m}$
$\longrightarrow A_{\bar m\leftarrow m}\,\operatorname{diag}(g)V_m$
$\longrightarrow$ 更新模板 $\longrightarrow$ 后续骨干与 CENTER 定位头
\end{minipage}}
\caption{搜索上下文条件化模板 Value 交互。新增检索分支读取本帧 early context，仅通过残差增益调制跨模态来源 Value；同模态消息沿原路径计算。}
\label{fig:framework}
\end{figure*}

\subsection{模板 MHSA 的消息分解}
在融合层 $\ell$，将当前两模态模板拼接为 $Z_\ell=[Z_{\ell,\R};Z_{\ell,\T}]$。
LayerNorm 和线性投影产生每个头的 $Q_\ell,K_\ell,V_\ell$：
\begin{equation}
 A_\ell=\softmax_{\mathrm{key}}\!\left(\frac{Q_\ell K_\ell^{\mathsf T}}{\sqrt{d_h}}\right).
 \label{eq:attention}
\end{equation}
softmax 在拼接后的全部模板 key 上进行，不对 self/cross 子块分别归一化。
将 $A_\ell$ 按接收模态和来源模态划分，可将接收模态 $n$ 的原生消息写为
\begin{equation}
 U_{\ell,n}=A_{\ell,n\leftarrow n}V_{\ell,n}
              +A_{\ell,n\leftarrow m}V_{\ell,m},\quad m\ne n.
\end{equation}
本文仅调制第二项。
多头输出拼接、输出投影及模板残差连接仍沿用原有定义。

\subsection{原有 early evidence}
对归一化后的 early token，定义同模态模板到搜索的相似度
\begin{equation}
 s_{jk}^{m}=\left\langle\frac{Z_{m,j}^{e}}{\|Z_{m,j}^{e}\|_2},
                              \frac{X_{m,k}^{e}}{\|X_{m,k}^{e}\|_2}\right\rangle.
\end{equation}
归一化在实现中使用 $10^{-6}$ 数值保护。
同一空间索引 $j$ 的原有证据为
\begin{equation}
 e_j=\left[\max_k s_{jk}^{\R},\ \max_k s_{jk}^{\T},\
             \cos(Z_{\R,j}^{e},Z_{\T,j}^{e})\right].
 \label{eq:evidence}
\end{equation}
原有预测器 $f_\ell:\mathbb{R}^3\to\mathbb{R}^2$ 使用隐藏宽度 32 的两层 MLP 和 GELU。
输出的两个坐标分别对应 RGB 来源与 TIR 来源。
三个输入是匹配摘要，不含显式前景标注，也不是定位效用标签。

\subsection{来源模板检索对侧搜索上下文}
为使增益预测器读取搜索内容，C 对两个传播方向使用共享的检索参数。
各融合层具有各自的检索模块。
对来源模态 $m$ 和模板位置 $j$，查询为
\begin{equation}
 q_{\ell,j}^{m}=W_{Q,\ell}\LN_Q(Z_{m,j}^{e}).
\end{equation}
将原有三维证据转换为三枚摘要 token：
\begin{equation}
 E_{\ell,j,a}=e_{j,a}b_{\ell,a},\quad a\in\{1,2,3\},
\end{equation}
其中 $b_{\ell,a}\in\mathbb{R}^{d}$ 是可学习基向量。
每个来源位置对应的记忆为
\begin{equation}
 \mathcal M_{\ell,j}^{m\to\bar m}
 =[E_{\ell,j,1};E_{\ell,j,2};E_{\ell,j,3};X_{\bar m}^{e}].
 \label{eq:memory}
\end{equation}
在记忆上施加共享 LayerNorm，再投影为 key 与 value。
每个头内的检索输出为
\begin{equation}
 c_{\ell,j}^{m\to\bar m}
 =\softmax\!\left(\frac{q_{\ell,j}^{m}K(\mathcal M)^{\mathsf T}}{\sqrt{r_h}}\right)V(\mathcal M).
 \label{eq:context}
\end{equation}
式~\eqref{eq:context} 为逐头表达，头输出拼接后再送入预测器。
当前实现的总投影宽度为 $r=64$，头数为 4，每头宽度 $r_h=16$。
三枚摘要与所有搜索 token 共同参与 softmax，而非分别归一化后等权相加。
注意力 softmax 使用 FP32，聚合权重转换为相应 value 精度。

拼接查询、检索内容与原有证据，生成方向对应的残差 logit：
\begin{equation}
 \delta_{\ell,j}^{m\to\bar m}
 =h_\ell([q_{\ell,j}^{m},c_{\ell,j}^{m\to\bar m},e_j]).
\end{equation}
$h_\ell$ 为 $131\to64\to1$ 的 MLP，使用 GELU。
这一分支保留被检索内容，补充式~\eqref{eq:evidence} 的三个标量。
检索相关性仍不能直接视为有益消息的判据，任务效用由训练及学习后干预检验。

\subsection{残差增益与跨模态 Value 调制}
将两个方向的残差分别加入原有预测器对应输出：
\begin{align}
 \eta_{\ell,j}^{m\to\bar m}
   &=f_\ell(e_j)_m+\delta_{\ell,j}^{m\to\bar m},\\
 \lambda_{\ell,j}^{m\to\bar m}
   &=\tanh(\alpha_\ell)\tanh(\eta_{\ell,j}^{m\to\bar m}),\\
 g_{\ell,j}^{m\to\bar m}
   &=\clip\!\left(\exp\{\lambda_{\ell,j}^{m\to\bar m}\},\,0.25,\,4\right).
 \label{eq:gain}
\end{align}
$\alpha_\ell$ 是可学习标量，初始值为 0.1。
接收模态 $n$ 的消息变为
\begin{align}
 \widetilde U_{\ell,n,i}
 &=\sum_j A_{\ell,n\leftarrow n,ij}V_{\ell,n,j}\nonumber\\
 &\quad+\sum_j A_{\ell,n\leftarrow m,ij}
           g_{\ell,j}^{m\to n}V_{\ell,m,j},\quad m\ne n.
 \label{eq:message}
\end{align}
增益沿来源列 $j$ 作用，对接收位置 $i$ 和注意力头共享。
它改变来源幅度，不改变单个来源 Value 的方向。
增益之后不重新归一化注意力，因此拒绝部分 cross 消息不会自动增加 self 消息预算。
当前层相同输入下的式~\eqref{eq:attention} 保持不变；后续层不享有这一不变性。

\subsection{初始化、训练与时间反馈边界}
新增残差预测器的末层在骨干外层初始化后重新置零。
因此，C 的初始残差为零，在相同公共参数下退化为原有 Early gate。
这里的恒等参照是原有门控路径，不要求其自身增益恒为 1。
首个实际训练输入的核验记录显示，零残差与关闭新增分支的三张定位输出图最大差为 0，保存加载也一致。
这验证了所测输入的实现一致性，不代表训练后的不同 checkpoint 输出相同。
模型从公共 SOT 预训练权重开始，新增参数随机初始化后按上述规则设置末层。
训练使用 fresh optimizer/scheduler，单阶段联合优化全部可训练参数，不由 M1 或 Early gate 权重续训。

C 及最终 Ours 不增加 utility loss，沿用分类热图、GIoU 与 $L_1$ 跟踪目标：
\begin{equation}
 \mathcal L=\mathcal L_{\mathrm{cls}}
               +2\mathcal L_{\mathrm{GIoU}}+5\mathcal L_{1}.
\end{equation}
表中的 utility-loss 候选是独立对照，不是 Ours 的组成部分。

首帧模板图像在每帧重新编码，融合后的模板 token 不跨帧写回。
模板更新及 temporal tokens 在当前评测中关闭。
跨帧反馈来自预测框决定下一帧搜索 crop。
因此，帧内消息传播与跨帧裁剪反馈需要分别分析；关闭模板更新不等于消除反馈漂移。

\section{实验与分析}
\subsection{数据、指标与比较角色}
主要评测覆盖 LasHeR 全部 244 条测试序列、219,292 帧。
报告 AUC/Success、Precision（PR）、Normalized Precision（NPR）、OP50 和 OP75，均以百分数表示。
差值使用百分点（pp）。
AUC 由序列 success curve 均值计算，再对序列等权汇总；它不同于直接对全部帧求平均 IoU。
PR 取 20 像素中心误差阈值，NPR 取归一化中心误差阈值 0.2。
OP50/OP75 按既有评测器在相应 overlap 阈值处读取。

B0 表示早期 canonical 基础系统，M1 表示已有模板 MHSA 系统，Ours 表示 C 加提前降 LR 的当前候选。
三者用于呈现系统发展及已有性能，不构成同环境下的单变量消融。
B0 与 M1 来自独立历史运行；M1 与新候选也有软件、采样和调度差异。
历史 Early gate P 单独标明来源风险，不与新配对控制混用。

新 C 运行使用 LasHeR train977 固定清单，每个采样 epoch 为 59,904 对，15 个采样 epoch 共 898,560 对。
batch 为 16，梯度累积为 8，有效 batch 为 128，对应 7,020 次成功 optimizer 更新。
这里的训练量是预先规定的完整训练预算，不意味着每个训练视频的全部帧逐一遍历。
短预算运行及未完成测试不进入主结果表。

\subsection{训练与推理设置}
C 和 Ours 从相同公共 SOT 权重独立开始单阶段训练，seed 为 42。
AdamW 初始学习率为 $10^{-4}$，骨干学习率乘以 0.1，weight decay 为 $10^{-4}$。
训练启用 AMP，梯度裁剪范数为 0.1；模板/搜索输入为 $128^2/256^2$，patch stride 为 16。
训练环境为 1$\times$ RTX 4090 D 24 GiB、Python 3.8.20、PyTorch 2.0.0+cu118、CUDA 11.8。
DataLoader 使用 4 workers、prefetch factor 2，persistent workers 关闭。
推理使用 FP32、串行 tracker，模板更新关闭。
M1 历史环境为相同卡型、Python 3.8.20、PyTorch 1.9.0+cu111、CUDA 11.1。
这些差异限制了相对 M1 的模块归因。

C 的原调度在更新 4,680 后将学习率乘以 0.1。
Ours 仅将这次衰减移至更新 936 后，终点仍为更新 7,020。
实际边界由运行入口明确设置。
两次训练的初态 hash 一致，前 7,488 个增强 microbatch 完全匹配，更新 936 后模型 state hash 也一致。
这一前缀核验支持将该配对解释为同结构的调度对照。
两者均评测预定终点权重；提前衰减的设计使用了已查看的开发与 full244 结果，因此仍需独立确认。

\subsection{核心系统性能}
表~\ref{tab:main} 给出 B0、M1 和 Ours 的完整测试结果。
Ours 相对 B0 的五项点估计均超过 $+1.5$ pp。
对更接近当前结构的 M1，增量主要体现在 OP75，而非 PR 的同等幅度改善。
表~\ref{tab:delta} 明确区分这两个参照，避免将基础系统差值写成相对 M1 的机制收益。

\begin{table*}[t]
\centering
\caption{核心系统在 LasHeR full244 上的已有结果。三行均完成各自预定完整训练与测试，但训练环境及历史条件未严格配对。B0 为本研究的 canonical 基础系统记录；Ours 为 C 结构加提前降 LR。粗体仅比较本表三行。}
\label{tab:main}
\small
\setlength{\tabcolsep}{3pt}
\begin{tabular}{@{}llrrrrrr@{}}
\toprule
系统 & 训练进度 & AUC$\uparrow$ & PR$\uparrow$ & NPR$\uparrow$ & OP50$\uparrow$ & OP75$\uparrow$ & FPS\\
\midrule
B0：TBSI 基础系统 & 15 epoch & 54.82 & 68.53 & 64.59 & 66.52 & 45.80 & NA\\
M1：模板 MHSA & 15 epoch & 56.21 & 70.21 & 66.20 & 68.33 & 47.13 & NA\\
Ours：C + 提前降 LR & 7,020 updates & \textbf{56.8305} & \textbf{70.4502} & \textbf{66.8347} & \textbf{68.8178} & \textbf{49.0406} & 25.58\\
\bottomrule
\end{tabular}
\end{table*}

\begin{table}[t]
\centering
\caption{Ours 相对两个历史参照的点估计差值（pp）。M1 差值用原归档精确值计算，再四舍五入，故不必等于主表显示值的直接相减。差值不表示单变量因果效应。}
\label{tab:delta}
\small
\setlength{\tabcolsep}{3pt}
\begin{tabular}{@{}lrrrrr@{}}
\toprule
参照 & AUC & PR & NPR & OP50 & OP75\\
\midrule
B0 & +2.01 & +1.92 & +2.24 & +2.30 & +3.24\\
M1 & +0.62 & +0.24 & +0.64 & +0.48 & +1.91\\
\bottomrule
\end{tabular}
\end{table}

\subsection{同结构学习率调度对照}
表~\ref{tab:schedule} 比较 C 原调度与提前衰减。
提前衰减的 AUC 和 OP75 点估计分别提高 0.2225 和 0.5837 pp。
AUC 逐序列胜/平/负为 130/3/111，配对序列 bootstrap 的 95\% 区间为 $[-0.7402,1.2625]$ pp。
OP75 的区间为 $[-0.2764,1.5830]$ pp。
这些区间均跨零，单 seed 结果不足以建立稳定调度优势。
由于结构未改变，这组差值也不构成新增 MHSA 机制的证据。

\begin{table*}[t]
\centering
\caption{C 架构的单次学习率衰减边界对照。两次运行均从公共预训练权重独立开始，终点为 7,020 次更新。区间仅度量给定 checkpoint 的序列抽样不确定性。}
\label{tab:schedule}
\small
\setlength{\tabcolsep}{3pt}
\begin{tabular}{@{}lrrrrrr@{}}
\toprule
设置 & AUC & PR & NPR & OP50 & OP75 & FPS\\
\midrule
C：更新 4,680 后衰减 & 56.6081 & 70.2697 & 66.5861 & 68.6029 & 48.4569 & 25.89\\
Ours：更新 936 后衰减 & 56.8305 & 70.4502 & 66.8347 & 68.8178 & 49.0406 & 25.58\\
差值（pp） & +0.2225 & +0.1805 & +0.2486 & +0.2150 & +0.5837 & --\\
95\% 序列区间 & $[-0.7402,1.2625]$ & $[-1.0601,1.4956]$ & $[-0.9130,1.5071]$ & $[-0.9963,1.5118]$ & $[-0.2764,1.5830]$ & --\\
\bottomrule
\end{tabular}
\end{table*}

\subsection{搜索条件表示与上下文层级}
为区分额外预测容量与搜索内容，已有短预算实验构建 P/S/C 三臂。
P 关闭新增残差；S 使用相同参数布局，但记忆仅含三枚证据摘要；C 增加对侧搜索 token。
各臂由公共预训练权重配对训练 936 次更新，在固定 screen96 开发 cohort 上评测。
该 cohort 限制每序列最多 2,000 帧，不是独立 full244 验证。
表~\ref{tab:short} 仅支持该开发范围内的比较。

\begin{table}[t]
\centering
\caption{936-update 的 P/S/C 开发对照，评测为 screen96。数值来自同一配对报告，不与 full244 指标混排。}
\label{tab:short}
\footnotesize
\setlength{\tabcolsep}{2pt}
\begin{tabular}{@{}lrrrrr@{}}
\toprule
设置 & AUC & PR & NPR & OP50 & OP75\\
\midrule
P：原有增益 & 55.5806 & 68.2487 & 65.0100 & 67.4259 & 47.1735\\
S：摘要记忆 & 54.5319 & 66.8962 & 63.6836 & 66.0791 & 46.1928\\
C：搜索记忆 & 56.4327 & 69.1302 & 65.9749 & 68.4986 & 47.9274\\
\bottomrule
\end{tabular}
\end{table}

C$-$P 的 AUC 为 $+0.8520$ pp，95\% 序列区间为 $[-0.4221,2.4861]$ pp。
C$-$S 为 $+1.9007$ pp，区间为 $[-0.1039,4.4025]$ pp；S$-$P 则为 $-1.0487$ pp。
因此，额外摘要处理容量本身未带来正向点估计；显式搜索内容在当前短预算下提供了正向信号。
这些结果不能替代相同 full 预算和调度下的 P/S/C 重训对照。
已有分析还显示，短预算的开发收益并未在后续 full 终点的同 cohort 上保持。
相关 full P 参照的权重及预测已不完整保留，故该变化也不能单独归因于训练时长。

另一完整预算实验仅将 C 新增检索分支的输入，从固定 early context 改为各融合层当前输入。
原有 early evidence、Value 动作、损失及 u936 衰减均保持不变。
其 AUC 为 56.4307，相对 Ours 下降 0.3998 pp；其余四指标也下降。
这暂不支持当前 layer-aligned 输入的有效增量，不能据此排除所有层级条件化设计。

\subsection{逐序列分布与历史来源敏感性}
相对 M1 的原归档指标，Ours 在 full244 上的 AUC 为 142 胜、0 平、102 负。
平均差为 $+0.6228$ pp，中位数为 $+0.2890$ pp。
配对序列 bootstrap 的 95\% 区间为 $[-0.5951,1.8906]$ pp。
该区间采用 20,000 次配对序列重采样，随机种子为 20261101，不校正多轮测试集选择。
有 23 条序列退化超过 5 pp，15 条超过 10 pp。
五个最大正增量占全部正增量的 41.6\%；去掉三个最大增益后，剩余序列均值为 $-0.0977$ pp。
这些是事后敏感性描述，不能据此重定义正式测试集。

M1 的七条现存预测与历史归档存在已知来源漂移。
正式 full244 仍保留原归档指标，不以现存目录重算值替换。
表~\ref{tab:sensitivity} 补充其余 237 条标量一致序列，以及其中由 M1 历史行为选出的 Q154。
标量一致不证明历史逐帧轨迹完全相同；Q154 也不是独立验证集。
这两组的 AUC 差接近零或为负，说明 full244 历史差值对来源与序列构成敏感。
七条风险序列的贡献不能直接判为虚假，也不能忽略。

\begin{table}[t]
\centering
\caption{Ours 相对 M1 的来源敏感性分析（pp）。237 条为历史标量一致范围；Q154 由 M1 行为事后选择。所有区间均为固定 checkpoint 的序列 bootstrap。}
\label{tab:sensitivity}
\footnotesize
\setlength{\tabcolsep}{2pt}
\begin{tabular}{@{}lrrr@{}}
\toprule
范围 & $\Delta$AUC & $\Delta$OP75 & AUC 95\% 区间\\
\midrule
完整 244 条 & +0.6228 & +1.9057 & $[-0.5951,1.8906]$\\
标量一致 237 条 & +0.0376 & +1.2843 & $[-0.9614,1.0832]$\\
事后 Q154 & -0.1640 & +1.2264 & $[-1.6482,1.3812]$\\
\bottomrule
\end{tabular}
\end{table}

在 237 条范围中，IoU$<0.2$ 的帧数由 M1 的 62,984 增至候选的 63,544。
按既有窗口事件定义，稳定后崩溃的序列数由 154 减至 150。
两个读数方向不同，不能用事件数下降概括为普遍恢复改善。
自由运行的两方法使用不同 crop，逐帧关联也不是相同输入下的反事实。
官方属性分层同样呈现异质性：该范围内低照度与高照度属性的 AUC 差为 $-2.5251/-1.7571$ pp。
这些探索性读数不支持普遍适应模态退化的主张。

\subsection{最终 checkpoint 的固定 crop 动作诊断}
学习后的动作是否即时改善定位，需要在相同 state 与 crop 下检验。
已有诊断加载 Ours 的 u7020 权重，在八条已查看的尾部序列第 50 帧比较原生增益和 $g=1$。
样本包含五条大退化序列与三条大增益序列，属于结果导向选样的机制探针。
GT 仅用于推理后评分。
八组比较的 state/crop hash 均一致，原生整数输出框与保存的 full244 输出一致。

$g=1$ 通过同时归零原有 Early gate 及搜索残差的 logit 实现。
因此，它检验最终来源增益的整体执行作用，不能独立归因搜索分支。
原生来源 log-gain 在这些 crop 和三个融合层中约处于 $[-0.0357,0.0808]$。
虽然动作非恒等，原生相对 $g=1$ 的平均有符号 IoU 差仅为 $-0.000014$，平均绝对差为 $0.000112$。
这八个固定状态暂不支持可靠的即时定位收益。
它们也不能证明其他状态上的动作无效，或排除该路径在训练过程中的影响。
该结果尤其不能外推为闭环恢复或错误传播已被纠正。

\subsection{效率、限制与结论边界}
Ours 的已记录吞吐为 25.58 FPS，C 原调度为 25.89 FPS。
吞吐由本次记录的帧数和跟踪耗时汇总，不扩展为未经单独测量的纯网络延迟或端到端部署速度。
历史 P 的 53.14 FPS 与其他候选速度使用各自运行条件，不能直接解释为纯结构开销。
相同 screen96 合同下，P 与 C 的 FPS 为 30.97 和 26.17，表明内容检索存在实际计算成本。
参数增量及统一合同下的显存、延迟对照尚不完整，不以估算值填补。

当前证据支持三个有边界的判断。
首先，搜索内容已被显式接入来源 Value 增益预测，并完成规定预算的训练与 full244 测试。
其次，Ours 获得较高的整体性能点估计，且同结构配对显示调度可以改变最终结果。
最后，来源风险、尾部抵消及固定 crop 的微小动作效应，限制了机制解释。
当前尚未区分结构、训练条件与历史比较差异对 M1 差值的相对贡献。

full244 已参与多轮候选和调度选择，不是完全未接触的最终确认集。
本文核心组只提供 seed 42，不以序列 bootstrap 代替跨 seed 方差。
也未建立独立数据集上的泛化结论或对最新外部方法的统一优势。
要确认条件化本身的收益，仍需保留完整产物的同调度 P/S/C 配对训练，以及独立 seed 或数据上的确认。
这些证据缺口不会改变已报告的性能记录，但决定了当前论文主张的范围。

\section{结论}
本文研究当前搜索内容如何参与模板跨模态 Value 交互。
所构建的 C 以来源模板检索对侧搜索内容，生成 Early gate 的残差增益。
采用提前学习率衰减的候选在 LasHeR full244 上取得 AUC 56.8305 和 OP75 49.0406。
相对 B0 的五项点估计均超过 $+1.5$ pp；相对 M1 的 AUC/OP75 为 $+0.62/+1.91$ pp。
同结构配对将调度贡献与结构主张分开，逐序列与固定 crop 诊断则揭示当前解释边界。
现有结果支持将该结构作为经验性能候选，尚不支持稳健的跨 seed 增益、退化识别或错误传播纠正结论。

\appendices
\section{完成完整预算与 full244 测试的候选}
表~\ref{tab:all} 保留核心候选组的全部已完成结果。
这些行来自不同方法和运行合同，不是逐行累加的组件消融。
特别是历史 P 的 PR 70.58 高于 Ours 的 70.4502，不能声称 Ours 在全部已列方法的五项指标上均为最高。
Ours 在该组的 AUC、NPR、OP50、OP75 点估计最高。

\begin{table*}[t]
\centering
\caption{完成各自完整训练预算与 LasHeR full244 测试的核心组。指标为百分数；NA 表示无已记录值。历史 P 有预测来源风险；速度未统一为纯网络计时。}
\label{tab:all}
\small
\setlength{\tabcolsep}{3pt}
\begin{tabular}{@{}llrrrrrr@{}}
\toprule
方法/实验 & 训练进度 & AUC & PR & NPR & OP50 & OP75 & FPS\\
\midrule
TBSI canonical B0 & 15 epoch & 54.82 & 68.53 & 64.59 & 66.52 & 45.80 & NA\\
M1 / TBSI-E-MHSA & 15 epoch & 56.21 & 70.21 & 66.20 & 68.33 & 47.13 & NA\\
Early gate 历史 P & 15 epoch & 56.56 & \textbf{70.58} & 66.76 & 68.69 & 47.64 & 53.14\\
EVDR aligned & 7,020 updates & 56.0956 & 69.6379 & 66.0366 & 67.8290 & 48.1135 & 58.72\\
$R_{\mathrm{sep}}$ & 7,020 updates & 56.4061 & 70.1943 & 66.2966 & 68.2390 & 48.2813 & 29.27\\
Early gate P + utility loss & 7,020 updates & 56.5590 & 70.3386 & 66.5281 & 68.4519 & 48.2914 & 31.77\\
Search-context C & 7,020 updates & 56.6081 & 70.2697 & 66.5861 & 68.6029 & 48.4569 & 25.89\\
Ours：C + 提前降 LR & 7,020 updates & \textbf{56.8305} & 70.4502 & \textbf{66.8347} & \textbf{68.8178} & \textbf{49.0406} & 25.58\\
C layer-aligned context & 7,020 updates & 56.4307 & 69.9619 & 66.4300 & 68.3135 & 48.5553 & 23.83\\
\bottomrule
\end{tabular}
\end{table*}

\section{短预算完整测试与未完成结果}
表~\ref{tab:partial} 区分训练预算和测试覆盖率。
C u936 与 T936 已完成 full244，但不是完整训练预算的结果。
T u7020 虽完成训练，现有记录仅有 62/244 条测试，故不提供可冒充 full244 的总体指标。
这里的 T 是独立候选标签，不是本文模态符号 $\T$ 所指的热红外分支。

\begin{table*}[t]
\centering
\caption{短预算完整测试与未完成测试，均不进入核心系统主表。}
\label{tab:partial}
\small
\setlength{\tabcolsep}{3pt}
\begin{tabular}{@{}llrrrrrl@{}}
\toprule
运行 & 训练/测试范围 & AUC & PR & NPR & OP50 & OP75 & 说明\\
\midrule
C u936 & 936 updates / full244 & 56.7486 & 70.3288 & 66.8916 & 68.9671 & 48.6825 & 短预算权重\\
T936 & 936 updates / full244 & 56.3947 & 69.9338 & 66.4862 & 68.5188 & 48.3948 & 短预算权重\\
T u7020 & 7,020 updates / 62 of 244 & NA & NA & NA & NA & NA & 测试未完成\\
\bottomrule
\end{tabular}
\end{table*}

\balance
\bibliographystyle{IEEEtran}
\bibliography{references}
\end{document}
